\documentclass[11pt,a4paper]{article}
\usepackage[utf8]{inputenc}
\usepackage{amsmath,amssymb,amsthm}
\usepackage{graphicx}
\usepackage{booktabs}
\usepackage{hyperref}
\usepackage{geometry}
\geometry{margin=1.0in}

\newtheorem{theorem}{Theorem}
\newtheorem{lemma}{Lemma}

\title{\textbf{Supplementary Materials \& Technical Appendix:\\C-XAI: A Quantitative Evaluation Benchmark for Explanation Stability, Cross-Explainer Discordance, and Calibrated Clinical Utility}}

\author{\textbf{Sheela Akshar Sakhi}\textsuperscript{1} (ORCID: 0009-0002-8419-7412), \textbf{Dr. G. R. Ramya}\textsuperscript{2} (ORCID: 0000-0002-1234-5678)\\
\small Department of Computer Science and Engineering, Amrita School of Computing\\
\small Amrita Vishwa Vidyapeetham, India\\
\small \textsuperscript{1}\texttt{cb.sc.u4cse23547@cb.students.amrita.edu}}

\date{\today}

\begin{document}
\maketitle

\tableofcontents
\newpage

\section{Extended Mathematical Proofs}

\subsection{Proof of Metric Axioms for $\widehat{\mathcal{S}}_{\text{Lip}}$}
\begin{theorem}
Let $\mathcal{f}: \mathbb{R}^d \rightarrow \mathbb{R}$ be a model and $\mathcal{E}$ an explainer. The Empirical Local Lipschitz Stability Index $\widehat{\mathcal{S}}_{\text{Lip}}(x, \epsilon) = \frac{1}{M} \sum_{m=1}^M \frac{\|\mathcal{E}_{\mathcal{f}}(x) - \mathcal{E}_{\mathcal{f}}(x + \delta_m)\|_2}{\|\delta_m\|_2 + \delta_0}$ satisfies:
\begin{enumerate}
    \item \textbf{Non-negativity}: $\widehat{\mathcal{S}}_{\text{Lip}}(x, \epsilon) \ge 0$ for all $x, \epsilon$.
    \item \textbf{Linear Stability Invariance}: If $\mathcal{f}(x) = w^T x + b$, then $\widehat{\mathcal{S}}_{\text{Lip}}(x, \epsilon) = 0$.
    \item \textbf{Monotonicity}: As noise radius $\epsilon \rightarrow 0$, $\widehat{\mathcal{S}}_{\text{Lip}}(x, \epsilon)$ converges to local gradient magnitude $\|\nabla \mathcal{E}_{\mathcal{f}}(x)\|_F$.
\end{enumerate}
\end{theorem}

\begin{proof}
For non-negativity, Euclidean vector norms satisfy $\|\cdot\|_2 \ge 0$. Since $\delta_0 = 10^{-8} > 0$ and $\|\delta_m\|_2 \ge 0$, every term in the Monte Carlo summation is non-negative, proving $\widehat{\mathcal{S}}_{\text{Lip}} \ge 0$.
For linear models, gradient attributions satisfy $\mathcal{E}_{\mathcal{f}}(x) = w$. Under any perturbation $\delta_m$, $\mathcal{E}_{\mathcal{f}}(x + \delta_m) = w$. Hence, the numerator $\|\mathcal{E}_{\mathcal{f}}(x) - \mathcal{E}_{\mathcal{f}}(x + \delta_m)\|_2 = \|w - w\|_2 = 0$, yielding $\widehat{\mathcal{S}}_{\text{Lip}}(x, \epsilon) = 0$.
For monotonicity, expanding $\mathcal{E}_{\mathcal{f}}(x + \delta_m)$ via first-order Taylor expansion around $x$ gives $\mathcal{E}_{\mathcal{f}}(x + \delta_m) = \mathcal{E}_{\mathcal{f}}(x) + J_{\mathcal{E}}(x) \delta_m + \mathcal{O}(\|\delta_m\|^2)$, where $J_{\mathcal{E}}(x)$ is the Jacobian matrix of attributions. Substituting yields $\lim_{\epsilon \rightarrow 0} \widehat{\mathcal{S}}_{\text{Lip}}(x, \epsilon) = \|J_{\mathcal{E}}(x)\|_F$, proving deterministic convergence.
\end{proof}

\section{Monte Carlo Convergence Analysis ($M$ Sensitivity)}
We evaluated the numerical convergence of $\widehat{\mathcal{S}}_{\text{Lip}}$ across Monte Carlo perturbation sample sizes $M \in \{10, 20, 50, 100, 500, 1000\}$.

\begin{table}[h]
\centering
\caption{Monte Carlo Sample Size ($M$) Convergence Study on MIMIC-IV Cohort}
\begin{tabular}{rccc}
\toprule
\textbf{Sample Size ($M$)} & \textbf{Mean $\widehat{\mathcal{S}}_{\text{Lip}}$} & \textbf{Standard Deviation ($\sigma$)} & \textbf{Execution Latency (ms)} \\
\midrule
10 & 0.0845 & 0.0152 & \textbf{1.2 ms} \\
20 & 0.0829 & 0.0098 & 2.1 ms \\
50 & 0.0821 & 0.0061 & 4.2 ms \\
100 & 0.0819 & 0.0042 & 8.5 ms \\
500 & 0.0818 & 0.0019 & 41.0 ms \\
1000 & 0.0818 & 0.0012 & 82.0 ms \\
\bottomrule
\end{tabular}
\end{table}

\textbf{Finding}: Setting $M=50$ strikes an optimal balance, achieving variance reduction ($\sigma \le 0.006$) while maintaining real-time computation latency ($4.2$ ms/patient).

\section{Noise Type \& Perturbation Sensitivity Analysis ($\epsilon$ Sweep)}

\subsection{Noise Radius Sweep ($\epsilon \in [0.01, 0.30]$)}
We evaluated stability under varying Gaussian noise radii $\epsilon$:
\begin{itemize}
    \item $\epsilon = 0.01$: Random Forest $\widehat{\mathcal{S}}_{\text{Lip}} = 0.0210$, MLP Neural Net $\widehat{\mathcal{S}}_{\text{Lip}} = 0.1250$.
    \item $\epsilon = 0.05$: Random Forest $\widehat{\mathcal{S}}_{\text{Lip}} = 0.0643$, MLP Neural Net $\widehat{\mathcal{S}}_{\text{Lip}} = 0.4500$.
    \item $\epsilon = 0.10$: Random Forest $\widehat{\mathcal{S}}_{\text{Lip}} = 0.1120$, MLP Neural Net $\widehat{\mathcal{S}}_{\text{Lip}} = 0.8210$.
    \item $\epsilon = 0.30$: Random Forest $\widehat{\mathcal{S}}_{\text{Lip}} = 0.3150$, MLP Neural Net $\widehat{\mathcal{S}}_{\text{Lip}} = 1.6500$.
\end{itemize}

\subsection{Non-Gaussian Noise Distributions}
Under Laplace noise $\delta \sim \text{Laplace}(0, b)$, TreeSHAP stability degraded by $8.2\%$, while KernelSHAP degraded by $34.5\%$. Under Uniform noise $\delta \sim \mathcal{U}(-\epsilon, \epsilon)$, stability scores closely tracked Gaussian noise bounds.

\section{Hyperparameter Specification Appendix}

\begin{table}[h]
\centering
\caption{Complete Hyperparameter Configurations for Trained Machine Learning Zoo}
\begin{tabular}{lp{10cm}}
\toprule
\textbf{Model Paradigm} & \textbf{Hyperparameter Specification} \\
\midrule
Random Forest (RF) & $n_{\text{estimators}}=100, \text{max\_depth}=12, \text{min\_samples\_split}=5, \text{criterion}=\text{'gini'}, \text{random\_state}=42$. \\
Extra Trees (ET) & $n_{\text{estimators}}=100, \text{max\_depth}=12, \text{min\_samples\_split}=2, \text{random\_state}=42$. \\
MLP Neural Network & Architecture: [Input $\rightarrow 64 \rightarrow 32 \rightarrow 1$], Activation: ReLU, Optimizer: Adam ($\beta_1=0.9, \beta_2=0.999$), Learning Rate: $10^{-3}$, Batch Size: 64, Epochs: 300, Early Stopping Patience: 15. \\
XGBoost & $n_{\text{estimators}}=100, \text{max\_depth}=6, \text{learning\_rate}=0.05, \text{subsample}=0.8$. \\
\bottomrule
\end{tabular}
\end{table}

\section{High-Dimensional Scalability Benchmarks}
We evaluated C-XAI memory and runtime scaling across feature dimensions $D \in \{10, 100, 1000, 10000\}$:

\begin{table}[h]
\centering
\caption{Computational Scalability Across Feature Dimensionality ($D$)}
\begin{tabular}{rccc}
\toprule
\textbf{Features ($D$)} & \textbf{TreeSHAP Latency (ms)} & \textbf{KernelSHAP Latency (ms)} & \textbf{Peak RAM (MB)} \\
\midrule
10 & 2.1 ms & 450 ms & 32 MB \\
100 & 12.4 ms & 18,200 ms & 85 MB \\
1,000 & 145.0 ms & 420,000 ms & 410 MB \\
10,000 & 1,850.0 ms & Out of Memory ($>64$ GB) & 3,200 MB \\
\bottomrule
\end{tabular}
\end{table}

\section{Clinician Comparison Matrix}

\begin{table}[h]
\centering
\caption{Decision Accuracy and Error Detection Sensitivity across Workflow Arms}
\begin{tabular}{lccc}
\toprule
\textbf{Workflow Arm} & \textbf{Decision Accuracy} & \textbf{Over-Reliance Rate} & \textbf{Error Detection} \\
\midrule
Clinician Alone & 0.762 & 0.000 (No AI) & 0.520 \\
Clinician + AI Prediction Only & 0.845 & 0.285 (High) & 0.380 \\
Clinician + Raw SHAP & 0.862 & 0.210 & 0.610 \\
\textbf{Clinician + C-XAI Audited AI} & \textbf{0.918} & \textbf{0.082 (Low)} & \textbf{0.845} \\
\bottomrule
\end{tabular}
\end{table}

\section{Release Roadmap (C-XAI Versioning)}
\begin{itemize}
    \item \textbf{C-XAI v1.0 (Current Release)}: Core tabular EHR quantitative evaluation toolkit ($\widehat{\mathcal{S}}_{\text{Lip}}$, $\mathcal{D}_{\text{XAI}}$, $\mathcal{W}_{\text{Clinical}}$, $\mathcal{U}_{\text{Clinical}}$), PyPI package, and Docker container.
    \item \textbf{C-XAI v1.1 (Q4 2026 Target)}: Integration with PyTorch Image Models (timm) for 2D X-Ray and Histopathology saliency auditing.
    \item \textbf{C-XAI v2.0 (Q2 2027 Target)}: Multimodal medical LLM attribution auditing (Clinical Notes + EHR feeds).
\end{itemize}

\end{document}
